\documentclass[10pt,a4paper]{amsart}
\usepackage[margin=19mm]{geometry}
\usepackage{amsmath,amssymb,mathtools}
\usepackage[scr=rsfs,cal=euler]{mathalfa}
\usepackage{xfrac}
\usepackage[unicode,hidelinks,bookmarks=false]{hyperref}
\newcommand{\ie}{i.e.\ }
\newcommand{\cf}{cf.\ }
\newcommand{\resp}{respectively\ }
\newcommand{\Subsection}{\subsection}
\providecommand{\nobreakdash}{\nobreak\hbox{-}}
\setlength{\emergencystretch}{2em}
\multlinegap0pt
\newcommand{\R}{\mathbb{R}}
\usepackage{amssymb}
\usepackage{tikz}
\usepackage{xcolor}
\usepackage{enumerate}
\usepackage{cancel}
\newtheorem{theorem}{Theorem}
\newtheorem{lemma}{Lemma}
\newtheorem{proposition}{Proposition}
\newcommand{\Br}{{\rm{Br}}}
\newcommand{\Id}{{\bf{1}}}
\newcommand{\ZS}{\mathbb{S}}
\newcommand{\ZT}{\mathbb{T}}
\newcommand{\be}{\beta}
\newcommand{\cC}{{\mathcal C}}
\newcommand{\cj}{\mathcal J}
\newcommand{\de}{\delta}
\newcommand{\e}{\varepsilon}
\newcommand{\ga}{\gamma}
\newcommand{\la}{\lambda}
\newcommand{\om}{\omega}
\newcommand{\supp}{{\rm supp}}
\title{Restriction and decoupling estimates for the hyperbolic paraboloid in $\R^3$}
\author{Ciprian Demeter, Shukun Wu}
\date{Source: arXiv:2505.09037v2 (CC BY 4.0)}
\begin{document}
\maketitle

\noindent\emph{Source and licence.} Restriction and decoupling estimates for the hyperbolic paraboloid in $\R^3$. Version arXiv:2505.09037v2 (CC BY 4.0), distributed by arXiv under the Creative Commons Attribution 4.0 International licence, http://creativecommons.org/licenses/by/4.0/, as recorded in the arXiv licence metadata for this exact version and shown on the abstract page https://arxiv.org/abs/2505.09037v2. Full licence notice: \texttt{licenses/arxiv-2505.09037-CC-BY-4.0.txt}.\par\smallskip

\noindent\textit{Editorial scope.} Counted unit: the article's proposition broad-prop (source0.tex lines 1408--1422). The article's complete original proof of it is delivered with it (source0.tex lines 1423--1655, one \texttt{proof} environment of 233 lines). No mathematics is written, repaired or re-derived: the statement and the proof are the article's own lines, in the article's own notation, and no retained line is substituted or re-flowed.  Retained uncounted as the source-local context the proof needs: lines 472--475; lines 1198--1201; lines 1233--1243; lines 1257--1266; lines 1269--1280; lines 1318--1321; lines 1325--1328; lines 1389--1395.  Nothing was omitted from any retained span, so the package carries no omission record.  No retained line contains a citation, so the wrapper carries no bibliography and no bibliography entry.  No retained line contains a figure environment, an image-inclusion command or drawing code, so no image is delivered and the package declares no asset dependency.  Editorial formatting only, in addition: an AMS \texttt{amsart} preamble that restates the article's own declarations and macros used by the retained lines, namely the environments \texttt{theorem}, \texttt{lemma}, \texttt{proposition} on separate counters, together with the standard packages those lines need.  The retained environments print in this document as Theorem 1, Lemma 1, Proposition 1 in order of appearance, whereas the article prints the same items with the numbering of its own full document.  The complete native source of the article as distributed in the arXiv e-print for this exact version is bundled unchanged as \texttt{sources/178087/source0.tex}.
\begin{theorem}(Bilinear refined decoupling)
\label{refined-dec-thm}
For all $\e>0$, we have	$C(R)\lesssim_\e R^\e$.
\end{theorem}

\begin{equation}
\label{fiourf8ur8gut8gu8tu}
\e_0=\e^{1000}.
\end{equation}

\begin{lemma}
	\label{wpt}
	The wave packet decomposition satisfies the following properties.
	\begin{enumerate}\item $Ef=\sum_{T\in\bar\ZT}Ef_{T}$.
		\item $|Ef_{T}(x)|\lesssim R^{-1000}$ when $x\in B_R\setminus T$.
		\item $\supp f_{T}\subset 3\theta$ when $T$ has direction $V(\theta)$.
		\item $\{V(\theta)\}_{\theta\in\Theta}$ are $\gtrsim R^{-1/2}$-separated.
		\item $\bar\ZT(\theta)$ is $R^{O(\e_0)}$-overlapping.
		\item $\|Ef_T\|_{L^p(w_{B_R})}\lesssim R^{2(\frac{1}{p}-\frac{1}{2})} \|Ef_T\|_{L^2(w_{B_R})}$ for all $T\in\bar\ZT $, where $w_{B_R}$ is a weight that is $\sim1$ on $B_R$ and decreases rapidly outside $B_R$. 
	\end{enumerate}    
\end{lemma}

\begin{lemma}
	\label{lem: l2}
	Let $X$ be a union of $R^{1/2}$-balls, and let $f=\sum_{T\in\ZT}f_T$ be a sum of wave packets.
	Suppose for each $T\in \ZT$ there is a shading $Y(T)\subset T$ by $R^{1/2}$-balls in $X$ such that the number of $R^{1/2}$-balls intersecting $Y(T)$ is $\lesssim \lambda R^{1/2}$.
	Then
	\begin{equation}
		\nonumber
		\int_{X}\big|\sum_{T\in\ZT}Ef_{T}\Id_{Y(T)}\big|^2\lesssim (\la R)\|f\|_2^2.
	\end{equation}
\end{lemma}

\begin{proposition}
	\label{kakeya-prop}
	Let $\de\in(0,1)$.
	Let $(L,Y)_\de$ be a collection of $\de$-separated lines together with an $(\e_1, \e_2)$-two-ends, $\la$-dense shading. 
	Define $E_L=\bigcup_{\ell\in L}Y(\ell)$.
	Suppose that every $\de$-ball on $\ZS^{2}$ contains $\leq m$ points from the direction set $\{V(\ell):\ell\in L\}$, where $V(\ell)$ is the direction of $\ell$.
	Take $\mu=\de^{-2\e_1}m\la^{-3/4}\de^{-1/2}$. Then there exists a set $E_\mu\subset E_L$ such that $\# L(x)\lessapprox \mu$ for all $x\in E_\mu$, and
	\begin{equation}
		\nonumber
		|E_L\setminus E_\mu|\leq \de^{\e_1}|E_L|.
	\end{equation}
\end{proposition}

	\begin{equation}
		\label{broad-triangle}
		\Br_A \{F^\tau\}(x)\leq \Br_{A_1} \{F^{\tau}_1\}(x)+\Br_{A_2}\{F^{\tau}_2\}(x)+\ldots+\Br_{A_N}\{F^{\tau}_N\}(x).
	\end{equation}

	\begin{equation}
		\label{broad-bilinear}
		\Br_A \{F^{\tau}\}(x)\leq \max_{\tau_1,\tau_2:\;K^{-1}-\text{transverse}}|F^{\tau_1}(x)F^{\tau_2}(x)|^{1/2}.
	\end{equation}

\begin{lemma}[Pigeonholing]
\label{lcomb}Consider a finite collection of numbers $I_Q,\;Q\in\mathcal Q$, with $I_Q\in [L,2L]$. Assume there is a finite set $\Lambda$ and numbers $A\le I_{Q,\lambda}\le B$ such that for each $Q\in\mathcal Q$
$$I_Q\le C\sum_{\lambda\in\Lambda}I_{Q,\lambda}.$$
Then there are $\lambda\in\Lambda$, $L'$, and $\mathcal Q''\subset \mathcal Q$ such that $I_{Q,\lambda}\in [L',2L']$ for each $Q\in\mathcal Q''$,
$$ (\log B/A)^{-1}(\#\Lambda)^{-1}\#\mathcal Q\le\#\mathcal Q'' $$ and
$$(C\log B/A)^{-1}(\#\Lambda)^{-2}\sum_{Q\in \mathcal Q}I_Q\lesssim \sum_{Q\in\mathcal Q''}I_{Q,\lambda}.$$
\end{lemma}

\begin{proposition}
	\label{broad-prop}
	Assume $\e\ll 10^{-3}$ is small enough.
	For $R\gg 1$, let $K=R^{\e^{10}}$.
	Then there exists $C_\e>0$ such that for all $R\geq1$,
	\begin{equation}
		\label{mixed-norm}
		\int_{B_R}|\Br_{A} Ef|^p\leq C_\e R^{2\e}\|f\|_2^2\;\sup_{\theta:R^{-1/2}\text{-square}}\|f_\theta\|_{L^2_{avg}(\theta)}^{p-2}
	\end{equation}
	for $p=22/7$ and $A\geq R^{\e^{20}}$.
	Here $\|f_\theta\|_{L^2_{avg}(\theta)}$ is defined as
	\begin{equation}
		\|f_\theta\|_{L^2_{avg}(\theta)}^2:=|\theta|^{-1}\|f_\theta\|_2^2.
	\end{equation}
\end{proposition}

\begin{proof}
	Throughout this argument, we let $p=22/7$. Fix $\e$.
	We use induction on  $r$ to prove that for all $r\in[R^{\e^2},R]$ and $A\geq r^{\e^{20}}$,
	\begin{equation}
		\label{mixed-norm-2}
		\int_{B_r}|\Br_{A} Ef|^p\leq C_\e R^\e r^\e\|f\|_2^2\sup_{\theta:r^{-1/2}\text{-square}}\|f_\theta\|_{L^2_{avg}(\theta)}^{p-2}.
	\end{equation}
	The base case is when $r=R^{\e^2}$, which is trivial via the use of elementary inequalities,  as $R^{\e}=r^{\e^{-1}}\geq r^{100}$. 
	We will see that the number $n$ of steps in this iteration is $\sim \log\e/\log(1-\e)=O_\e(1)$. Indeed, the sequence of radii is $$R^{\e^2},\;R^{\e^2/(1-\e^2)},\ldots, R^{\e^2/(1-\e^2)^n}\sim R.$$
	
	\smallskip

    Assume \eqref{mixed-norm-2} holds for $r=R^{\e^2/(1-\e^2)^{m-1}}$, $m\ge 1$. Fix $r=R^{\e^2/(1-\e^2)^m}$ and fix $B_r$. Partition the $r$-tubes $\bar\ZT=\ZT\sqcup\ZT_{small}$, where 
    $\ZT_{small}=\{T\in\bar\ZT:\|f_T\|_2\leq r^{-100}\|f\|_2\}$.
	An easy computation shows that $ \int_{B_r}|\Br_{A} E(\sum_{T\in\ZT_{small}}f_T)|^p\lesssim r^{-10}\|f\|_2^p$, which trivially yields \eqref{mixed-norm}. It remains to estimate $\int_{B_r}|\Br_{A} Ef'|^p,$
    where $f'=\sum_{T\in\ZT}f_T$.
    Next, partition $\ZT=\bigsqcup_{\ga,m}\ZT_{\ga,m}$, where $\ga,m\in[r^{-100},r^{10}]$ are dyadic numbers, such that 
	\begin{itemize}
		\item For all $T\in\ZT_{\ga,m}$, $\|f_T\|_2\sim \ga\|f\|_2$\;.
		\item For all $\theta$, either $\ZT_{\ga,m}(\theta)=\varnothing$, or $\#\ZT_{\ga,m}(\theta)\sim m$.
	\end{itemize}
	Since there are $O((\log r)^2)$ possible pairs of dyadic numbers $(\ga,m)$, by the triangle inequality \eqref{broad-triangle}, there exists a pair  $(\ga,m)$ and  $A_g\gtrapprox A$ such that, writing $g=\sum_{T\in\ZT_{\ga,m}}f_T$, we have 
	\begin{equation}
		\label{jfdjiofgjbiogjbigji}
		\int_{B_r}|\Br_{A} Ef'|^p\lessapprox \int_{B_r}|\Br_{A_g} Eg|^p.
	\end{equation}
	To ease notation, we let $\ZT_g=\ZT_{\ga,m}$.
    By dyadic pigeonholing, there exists a union $X$ of $r^{1/2}$-balls $Q$ such that 
	\begin{itemize}
		\item The values  $\int_{Q}|\Br_{A_g} Eg|^p$ are about the same for all $Q\subset X$.
		\item We have
		\begin{equation}
			\label{reduction-1}
			\int_{B_r}|\Br_{A_g} Eg|^p\lessapprox\int_{X}|\Br_{A_g} Eg|^p.
		\end{equation}
	\end{itemize}
	
	
	
	\medskip
	
	\noindent{\bf Step 1: Two-ends reduction.}
	
	Partition each $r$-tube $T\in\ZT_g$ into  tube segments $J$ of length $r^{1-\e^2}$. Let $\cj(T)$ be those segments that intersect $X$. 
	Then, partition  $\cj(T)=\bigcup_\la\cj_\la(T)$, $\la\in\Lambda$, where $\Lambda$ denotes the dyadic numbers in $[r^{-1/2},r^{-\e^2}]$, and $|J\cap X|\sim \la |T|$ for any $J\in\cj_\la(T)$. 
	Thus, 
	\begin{equation}
		\nonumber
		Eg=\sum_{\la}\sum_{T\in\ZT_g}Ef_{T}\sum_{J\in\cj_\la(T)}\Id_J.    
	\end{equation}
    Write $F_\lambda^\tau=\sum_{T\in\ZT_g\atop{\theta_T\subset\tau}}Ef_{T}\sum_{J\in\cj_\la(T)}\Id_J$, $\tau\in\cC_K$. Note that $Eg_\tau=\sum_{\lambda}F_\lambda^{\tau}$. The triangle  inequality \eqref{broad-triangle} together with the triangle inequality in $L^p$ followed by H\"older imply that, for some $A_1\gtrapprox A_g$, we have for each $Q\subset X$
    $$\int_{Q}|\Br_{A_g} Eg|^p\le \#(\Lambda)^{p-1}\sum_{\lambda}\int_{Q}|\Br_{A_1} \{F_\lambda^{\tau}\}|^p.$$

    We may assume all nonzero terms $I_{Q,\lambda}:=\int_{Q}|\Br_{A_1} \{F_\lambda^{\tau}\}|^p$ are in the interval $[r^{-100}(\gamma\|f\|_2)^p,r^{100}(\gamma\|f\|_2)^p]$.    
	Since $I_Q:=\int_{Q}|\Br_{A_g} Eg|^p$ are about the same for all $Q\subset X$, and since $\#\Lambda\lessapprox 1$, by Lemma \ref{lcomb} there is a $\la\in\Lambda$ and a set of $r^{1/2}$-balls $X_1\subset X$ such that 
	\begin{itemize}
		\item $|X_1|\gtrapprox|X|$.
		\item For each $r^{1/2}$-ball $Q\subset X_1$, 
        $\int_{Q}|\Br_{A_1} \{F_\lambda^{\tau}\}|^p
        $ has about the same value.
		\item We have
		\begin{equation}			
			\int_{X}|\Br_{A_g} Eg|^p\lessapprox\int_{X_1}|\Br_{A_1} \{F_\lambda^{\tau}\}|^p.
		\end{equation}
	\end{itemize}
	Consider the partition $\ZT_g=\bigcup_\be\ZT_{\be}$, where $\be\in[1,r^{\e^2}]$ is a dyadic number and $\#\cj_\la(T)\sim\be$ for all $T\in\ZT_{\be}$. 
	As a result, 
	\begin{equation}
		\nonumber
		\sum_{T\in\ZT_g}\sum_{J\in\cj_\la(T)}Ef_{T}\Id_J=\sum_\be\sum_{T\in\ZT_{\be}}\sum_{J\in\cj_\la(T)}Ef_{T}\Id_J.    
	\end{equation}Write $F_{\lambda,\beta}^\tau=\sum_{T\in\ZT_\beta\atop{\theta_T\subset\tau}}Ef_{T}\sum_{J\in\cj_\la(T)}\Id_J$, $\tau\in\cC_K$. Note that $F_\lambda^\tau=\sum_{\beta}F_{\lambda,\beta}^{\tau}$.
	Reasoning as in the previous step, using the triangle  inequality \eqref{broad-triangle} and  Lemma \ref{lcomb}, 
     we find  $\be$, $A_2\gtrapprox A_1$ and a set of $r^{1/2}$-balls $X_2$ such that 
	\begin{itemize}
		\item $|X_2|\gtrapprox|X_1|$.
		\item For each $r^{1/2}$-ball $Q\subset X_2$, $\int_{Q}|\Br_{A_2} \{F_{\lambda,\beta}^{\tau}\}|^p$ has about the same value.
		\item We have
		\begin{equation}
			\label{reduction-2}
			\int_{X_1}|\Br_{A_1} \{F_\lambda^{\tau}\}|^p\lessapprox\int_{X_2}|\Br_{A_2} \{F_{\lambda,\beta}^{\tau}\}|^p .
		\end{equation}
	\end{itemize}
	
	\smallskip
    
    It remains to analyze the last integral. We will distinguish two cases.
	Let $\{B_k\}$ be a finitely overlapping family of $r^{1-\e^2}$-balls that cover $B_r$. 
	\medskip
	
	\noindent{\bf Step 2: The non-two-ends scenario.} Assume $\be\leq r^{\e^4}$.
	
	For each $B_k$, define 
	\begin{equation}
		\nonumber
		g_{k}=\sum_{\substack{T\in\ZT_{\be} \text{ such that}\\ \exists J\in\cj_\la(T),\, J\cap B_k\not=\varnothing}} f_{T}.
	\end{equation}
    Note that on each $B_k$, by Lemma \ref{wpt}
    $$\big|\sum_{T\in\ZT_{\be}}Ef_T\sum_{J\in\cj_{\la}(T)}\Id_J\big|\sim |Eg_k|.$$
	Thus, we have
	\begin{align}
		\label{related}
		\int_{X_2\cap B_k}\big|\Br_{A_2}\{F_{\lambda,\beta}^\tau\}\big|^p&\sim\int_{X_2\cap B_k}\big|\Br_{A_2} Eg_k\big|^p.
	\end{align}
	Note that for each $T$, there are $\lesssim r^{\e^{4}}$ many $B_k$ such that $\exists J\in\cj_\la(T), J\cap B_k\not=\varnothing$.
	As a consequence, 
	\begin{equation}
		\label{l2-related}
		\sum_{k}\|g_k\|_2^2\lesssim r^{\e^{4}}\|g\|_2^2\lesssim r^{\e^4}\|f\|_2^2. 
	\end{equation}
	Since $A_2\gtrapprox A\geq r^{\e^{20}}$, we have (for $\e$ small enough) $A_2\geq r^{(1-\e^2)\e^{20}}$.
	Apply \eqref{mixed-norm-2} as an induction hypothesis on each $r^{1-\e^2}=R^{\e^2/(1-\e^2)^{m-1}}$-ball $B_k$ to get
	\begin{equation}
		\label{vpofigtug8utr8gurtg98u}
		\|\Br_{A_2} Eg_k\|_{L^p(B_k)}^p\leq C_\e R^\e r^{(1-\e^2)\e}\|g_k\|_2^2\;\sup_{\om}\|g_{k,\om}\|_{L_{avg}^2(\om)}^{p-2},
	\end{equation}
	where the sup is over  $ r^{(\e^2-1)/2}$-squares $\om$.
	By $L^2$-orthogonality,  
    \begin{equation}
    \label{rpojirutgut8}
    \sup_{\om}\|g_{k,\om}\|_{L^2_{avg}(\om)}^{p-2}\lesssim\sup_{\theta}\|f_{\theta}\|_{L^2_{avg}(\theta)}^{p-2}.
    \end{equation}
	Summing up over $B_k$, using \eqref{jfdjiofgjbiogjbigji}-\eqref{rpojirutgut8},  we find
	\begin{align}
		\nonumber
		\int_{B_r}|\Br_{A} Ef'|^p&\lessapprox \, \sum_{k}C_\e R^\e r^{(1-\e^2)\e} \|g_k\|_2^2\sup_{\om}\|g_{k,\om}\|_{L^2_{avg}(\om)}^{p-2}\\ \nonumber
		&\lesssim  r^{-\e^3+\e^4}C_\e R^\e r^{\e}\|f\|_2^2\sup_{\theta}\|f_{\theta}\|_{L^2_{avg}(\theta)}^{p-2}.
	\end{align}
	This proves \eqref{mixed-norm-2}. Note that  we  have not used yet either the information (gained via pigeonholing) regarding the subsets $X_i$ or the constant property relative to $Q$. These will be used in the next step, more precisely, in the derivation of \eqref{bdbhdbvhhvbfdhvbhfb}.  
	
	\medskip
	
	\noindent{\bf Step 3: The two-ends case.} Assume $\be\ge r^{\e^4}$.
    
	For each $T\in\ZT_{\be}$, consider the shading $Y(T)=\bigcup_{J\in\cj_\la(T)}(J\cap X)$.
	Then $Y$ is a rescaled $(\e^2, \e^{4})$-two-ends, $\la\be$-dense shading.
	
	
	\smallskip
	
	
	
	
	Define
	\begin{equation}
		\nonumber
		\mu=r^{2\e^{2}}m(\la\be)^{-3/4}r^{1/4}.
	\end{equation}
	We  apply Proposition \ref{kakeya-prop} to the $r^{-1}$-dilate of $(\ZT_\be, Y)$ (with $\delta=r^{-1/2}$) to obtain a set $X_3\subset X$ with
	 $|X\setminus X_3|\leq r^{-\e^2}|X|$ and ( recall $\e_0$ from \eqref{fiourf8ur8gut8gu8tu})
    \begin{equation}
		\label{multi-R-half}
		\sup_{Q\subset X_3}\#\{T\in\ZT_\be:Y(T)\cap Q\not=\varnothing\}\lessapprox r^{O(\e_0)}\mu.
	\end{equation}
	
	Since $|X_2|\gtrapprox |X|$, we know that  $|X_2\setminus X_3|\lessapprox r^{-\e^2}|X_2|$. 
	Denote by $X_4=X_2\cap X_3$, so we have $|X_4|\gtrapprox|X_2|$ and $X_4\subset X_2$. 
	Recall that $\int_{Q}|\Br_{A_2} \{F_{\lambda,\beta}^{\tau}\}|^p$
     are about the same for $Q\subset X_2$.
	Thus,  we have 
	\begin{equation}
		\label{bdbhdbvhhvbfdhvbhfb}\int_{X_2}|\Br_{A_2} \{F_{\lambda,\beta}^{\tau}\}|^p
		\lessapprox \int_{X_4}|\Br_{A_2} \{F_{\lambda,\beta}^{\tau}\}|^p.
	\end{equation}
    The change of the domain of integration from $X_2$ to $X_4$ is crucial, as it will give us access to the incidence estimate \eqref{multi-R-half}.
    
	Assuming $\e$ is small enough, we have $A_2\ge 2$, as $A_2\gtrapprox A\ge r^{\e^{20}}$.
	We invoke \eqref{broad-bilinear} and pigeonholing to find two $K^{-1}$-transverse $\tau_1,\tau_2\in\cC_K$ so that, denoting $\ZT_{\be}[\tau_j]=\bigcup_{\theta\subset\tau_j}\ZT_{\be}(\theta)$, we have
	\begin{align}
		\nonumber
        \int_{X_4}|\Br_{A_2} \{F_{\lambda,\beta}^{\tau}\}|^p
		\lesssim K^{O(1)}\int_{ X_4}\prod_{j=1,2}\Big|\sum_{T\in\ZT_{\be}[\tau_j]}\sum_{J\in\cj_{\la}(T)}Ef_{T}(x)\Id_J\Big|^{p/2}.
	\end{align}
	Recall that $\{B_k\}$ is a partition of $B_r$ into $r^{1-\e^2}$-balls. 
	For each $B_k$, let $$\ZT_{\be, k}[\tau_j]=\{T\in\ZT_{\be}[\tau_j]: \exists \;J\in\cj_{\la}(T),\; J\cap B_k\not=\varnothing\}.$$
	Therefore, using earlier inequalities we find
	\begin{align}
		\label{X-4}
		\int_{B_r}|\Br_{A} Ef'|^p&\lessapprox K^{O(1)} \sum_k\int_{X_4\cap B_k}\prod_{j=1,2}\Big|\sum_{T\in\ZT_{\be}[\tau_j]}\sum_{J\in\cj_{\la}(T)}Ef_{T}(x)\Id_J\Big|^{p/2}\\ \nonumber
		&\sim K^{O(1)}\sum_k\int_{X_4\cap B_k}\prod_{j=1,2}\Big|\sum_{T\in\ZT_{\be, k}[\tau_j]}Ef_{T}\Big|^{p/2}\\ 
        %\label{each-k}
        \nonumber
		&\lesssim  K^{O(1)} r^{10\e^2}\max_k\int_{X_4\cap B_k}\prod_{j=1,2}\Big|\sum_{T\in\ZT_{\be, k}[\tau_j]}Ef_{T}\Big|^{p/2}.
	\end{align}
    We derive two estimates for the right-hand side. 
    
	Notice that for each $Q\subset B_k$,
	\begin{equation}
		\nonumber
		\{T\in\ZT_{\be}:Y(T)\cap Q\not=\varnothing\}=\{T\in\ZT_{\be,k}:T\cap Q\not=\varnothing\}.
	\end{equation}
    When combined with \eqref{multi-R-half}, this shows that when $Q\subset X_4\cap B_k$
    \begin{equation}
    \label{multi-R-half4}
    \#\{T\in\ZT_{\be, k}:T\cap Q\not=\varnothing\}\lessapprox r^{O(\e_0)}\mu.
    \end{equation}
	At this point, we invoke Theorem \ref{refined-dec-thm} at scale $r$, using the set $X_4\cap B_k\subset B_r$ and  the bound \eqref{multi-R-half4}  to have
	\begin{equation}
		\label{jkijurgu98tug8tug89t}
		\int_{X_4\cap B_k}\prod_{j=1,2}\Big|\sum_{T\in\ZT_{\be, k}[\tau_j]}Ef_{T}\Big|^2\lessapprox K^{O(1)} r^{O(\e_0)} \mu\sum_{T\in\ZT_g}\big\|Ef_T\big\|_{L^4(w_{B_r})}^4.
	\end{equation}
	
	Recall that $\|f_T\|_{2}$ have comparable magnitude for all $T\in\ZT_g$, and that $\#\ZT_g(\theta)\sim m$ for all $\theta$ such that $\ZT_g(\theta)\not=\varnothing$.
	Thus, for each $\theta'$ we have
	\begin{align}
		\nonumber
		\sum_{T\in\ZT_g(\theta')}&\big\|Ef_T\big\|_{L^4(w_{B_r})}^4\lesssim r^{-2}\sum_{T\in\ZT_g(\theta')}\big\|Ef_T\big\|_{L^2(w_{B_r})}^4\lesssim\sum_{T\in\ZT_g(\theta')}\|f_T\|_2^4\\
		\label{after-decoupling}
		&\lesssim m^{-1} \Big(\sum_{T\in\ZT_g(\theta')}\big\|f_T\big\|_{2}^2\Big)^2\lesssim  (mr)^{-1}\|f_{\theta'}\|_2^2\;\sup_\theta\|f_\theta\|_{L^2_{avg}(\theta)}^2.
	\end{align} 
	We recall that $O(\e_0)\leq\e^2$, $K=R^{\e^{10}}\leq r^{\e^8}$, and $\be\ge 1$.
	Thus, summing up over all $\theta'$ in \eqref{after-decoupling} and plugging it back into \eqref{X-4},\eqref{jkijurgu98tug8tug89t}, we have
	\begin{align}
		\label{l4}
		\int_{X_4\cap B_k}\prod_{j=1,2}\Big|\sum_{T\in\ZT_{\be, k}[\tau_j]}Ef_{T}\Big|^2&\lessapprox r^{O(\e^2)}\mu (m r)^{-1}\|f\|_2^2\sup_\theta\|f_\theta\|_{L^2_{avg(\theta)}}^2\\  \nonumber
		&\lessapprox r^{O(\e^{2})}(\la r)^{-3/4}\|f\|_2^2\;\sup_\theta\|f_\theta\|_{L^2_{avg(\theta)}}^2.
	\end{align}
    This gives us a first estimate.
	
	
	Since $|T\cap ( X\cap B_k)|\lesssim \la  |T|$ for all $T\in\ZT_{\beta,k}$, by Cauchy-Schwarz and by Lemma~\ref{lem: l2} we get a second estimate  
	\begin{equation}
		\label{l2-1}
		\int_{X_4\cap B_k}\prod_{j=1,2}\Big|\sum_{T\in\ZT_{\be, k}[\tau_j]}Ef_{T}\Big|\lesssim (\la  r)\|f\|_2^2.
	\end{equation}
	Therefore, since $K=r^{O(\e^2)}$, by \eqref{X-4}, $\eqref{l4}^{4/7}\cdot\eqref{l2-1}^{3/7}$ gives when $p=22/7$,
	\begin{align}
		\nonumber
		\int_{B_r}|\Br_{A} Ef'|^p\lessapprox r^{O(\e^2)} \|f\|_2^2\;\sup_\theta\|f_\theta\|_{L^2_{avg(\theta)}}^{p-2}\leq C_\e R^{\e}r^\e\|f\|_2^2\;\sup_\theta\|f_\theta\|_{L^2_{avg(\theta)}}^{p-2}.
	\end{align}
	This proves \eqref{mixed-norm-2} and hence Proposition \ref{broad-prop}. \qedhere
	
	
\end{proof}

\end{document}
